\scanpage{194}
\mspage{II.42}
% archive sheet number: 193
% block: upper third, repeated 6.13 | text
6.13. Le cas le plus intéressant est celui où les $X_\lambda$ sont de type
\deleted{\unsure{de $X_\lambda$}} fini sur $\mathbb{Z}$, ou sur un corps
premier, \deleted{Alors} On sait que le foncteur
\inserted{$\unsure{\mathscr{X}|\mathcal{X}}\mathbin{-}X$}
\unsure{limite} des pro-objets dans la catégorie des schémas de type fini
sur $\mathbb{Z}$, dans celle des schémas, est pleinement fidèle, i.e.
$\mathscr{X}$ est déterminé par la connaissance de $X$.

% block: middle, example | text
\uline{Ex} $X$ de présentation finie sur un schéma affine, p.ex. sur un
corps. En particulier $X$ affine, p.ex. \unsure{$=$|$X$} spectre d'un
corps, p.ex. d'un corps alg. clos (sic)

% block: middle, arithmetic group name | text
Dans ce cas, $\Pi$ mérite d'être \unsure{appellé|appelé} le
``\inserted{pro}\uline{groupe fondamental arithmétique de $X$}'',
\inserted{notation $\Pi_1^{\mathrm{arith}}(X,\xi)$}
(en se rappelant que c'est un \deleted{groupe} pro-groupe topologique).
On a un isom. canonique
% block: lower middle, canonical isomorphism | display
\[
\pi_1(X,\xi)\xrightarrow{\sim}\varprojlim\Pi
=\varprojlim_\lambda\pi_1(X_\lambda,\xi),
\]
% block: lower third, limitation | text
mais on n'\unsure{oubliera} que la connaissance de $\pi_1(X,\xi)$ ne
détermine pas celle \deleted{de} \inserted{pro-groupe}
\deleted{\unsure{du}} $\Pi=\Pi_1(X,\xi)$.
Par exemple, si $X$ est le spectre d'un corps alg. clos,
$\pi_1(X,\xi)=e$, tandis que $\Pi$ n'est alors jamais trivial, comme
nous allons voir.

% block: foot, continuation | text
En vertu de 6.10, $\unsure{\pi^{\mathrm{arith}}|\Pi^{\mathrm{arith}}}(X,\xi)$
donne
